\section{Introduction}
\label{sec:introduction}

The primitive residual isolated in TPC-18--20 has an exact centered
divisor expansion on every finite target horizon.  If
\[
 a(u)=-\mu(u)\log u,
\]
then its common four-channel coefficient is
\begin{equation}
 b_R(u)=a(u)-\lambda'_R(u),
 \qquad \lambda'_R(u)=0\quad(u>R).
 \label{eq:intro-bR}
\end{equation}
Thus \(b_R(u)=a(u)\) beyond the finite Ramanujan cutoff.  TPC-21
removes the deterministic Ramanujan mean of the selected packet
\(u,v\le R\), and TPC-22 closes its centered discrepancy in explicit
parameter regions by exact-conductor rectangularization
\citep{WangTPC18,WangTPC19,WangTPC20,WangTPC21,WangTPC22}.

The support remaining after that result has four quadrants:
\begin{equation}
 \begin{array}{c|cc}
  &v\le R&v>R\\ \hline
 u\le R&\mathrm{SS}&\mathrm{SL}\\
 u>R&\mathrm{LS}&\mathrm{LL}.
 \end{array}
 \label{eq:intro-four-quadrants}
\end{equation}
Only SS was previously closed.  In particular, both-large terms exist:
one may take \(u=m_1j+h\) and \(v=m_2j+h\).  We therefore avoid the
ambiguous phrase ``one-sided large-divisor terms'' for the whole
complement.  We also distinguish support legs from the
\emph{single-indicator Fourier channels} completed in TPC-20; the
latter theorem did not estimate SL or LS in
\eqref{eq:intro-four-quadrants}.

\subsection{A dynamic cutoff}

Our first result moves the closed support itself.  Choose
\begin{equation}
 R<S\le X^{1-o(1)},
 \qquad \vartheta_{R,S}(u)=b_R(u)\one_{u\le S}.
 \label{eq:intro-dynamic-coefficient}
\end{equation}
This does not replace the finite model cutoff \(R\): it truncates the
\emph{residual} divisor coefficient at a later point.  The enlarged
packet \(u,v\le S\) contains SS, both mixed collars, and the nonempty
both-large square \((R,S]^2\).

Write
\[
 Q=X^{\beta+o(1)},\qquad J=X^{1-\beta+o(1)},
 \qquad S=X^{s+o(1)}.
\]
The TPC-22 character coordinates, exact-conductor large sieve, and
all-frequency induced-Gauss estimate remain valid with \(S\) in place
of \(R\).  The actual off-diagonal row-pair mask is essential: it gives
\[
 \sum_{g\le S}\sum_{m,n}|c_{m,n}^{(g)}|^2
 \ll Q^2X^\eps
\]
because \(g\mid m-n\), not because \(S\) is below the row length.  No
new hypothesis \(L>2S\) is required for the hybrid argument.  What
changes is the denominator
range \(q\le S^2\), the conductor range \(F_i\le S\), and the
two-scale completion condition \(HS\ll JX^{-\eta}\).

The resulting hybrid saving is
\begin{equation}
 \eta_S(\beta,s)=
 \min\left\{
 1-\frac{3\beta}{2},
 \frac{\beta+1-4s}{2},
 \frac{3-\beta-6s}{4}
 \right\}.
 \label{eq:intro-etaS}
\end{equation}
There is also an occupancy route when \(2s<1-\beta\), provided the
stronger source margin \(L>2S\) is imposed.  Consequently,
putting
\[
 s=\frac12-\delta+\rho
\]
gives a genuine strip beyond \(R\).  In the hybrid region it is enough
that
\begin{equation}
 0<\rho<\min\left\{
 \frac12+\delta-\beta,
 \frac{\beta+4\delta-1}{4},
 \frac{6\delta-\beta}{6}
 \right\},
 \label{eq:intro-rho-hybrid}
\end{equation}
while the occupancy region permits
\begin{equation}
 0<\rho<\frac{2\delta-\beta}{2}.
 \label{eq:intro-rho-occupancy}
\end{equation}
Every bound is strict and carries a fixed margin.  These formulas show
that every strict interior point of the TPC-22 hybrid closure region has
a positive collar.  The occupancy branch has the same conclusion only
when its additional source margin \(L>2S\) is available.  In either
case, this is a rigorous estimate for genuinely new large-divisor
coefficients under the inherited packet interfaces stated below.

\subsection{What the next support looks like}

Let \(n=mj+h\).  The part beyond \(S\), after a uniformly negligible
scalar centering tail is removed, is
\begin{equation}
 \cT_{m,S}(j)
 =\sum_{\substack{u\mid n\\u>S}}a(u)
 =\sum_{\substack{r\mid n\\r<n/S}}a(n/r).
 \label{eq:intro-complement}
\end{equation}
The second equality is exact and valid without assuming that \(n\) is
squarefree.  Since \(\mu(k)=\lambda(k)\mu^2(k)\), it also gives the
Liouville-cocycle identity
\begin{equation}
 \cT_{m,S}(j)
 =-\lambda(n)
 \sum_{\substack{r\mid n\\r<n/S}}
 \lambda(r)\mu^2(n/r)\log(n/r).
 \label{eq:intro-Liouville-cocycle}
\end{equation}
Thus divisor reflection turns geometric size into a target-dependent
parity phase.

More usefully for two rows, fix complement divisors \(r,s\).  Their two
divisibility conditions select either no orbit points or one class
\(j=j_0+[r,s]t\).  The quotient divisors are affine forms
\[
 U(t)=\frac{m_1(j_0+[r,s]t)+h}{r},\qquad
 V(t)=\frac{m_2(j_0+[r,s]t)+h}{s},
\]
and
\begin{equation}
 \det(U,V)=\frac{h(m_1-m_2)}{(r,s)}\ne0.
 \label{eq:intro-affine-determinant}
\end{equation}
The remaining LL coefficient is therefore
\(\mu(U(t))\mu(V(t))\log U(t)\log V(t)\).  This exact normal form
identifies the next analytic input: a quantitatively uniform two-point
affine M\"obius estimate for growing coefficients, summed over moving
complement fibers.

For each fixed fiber, the logarithmically averaged estimate is already
known: Tao's two-point Elliott theorem applies because
\eqref{eq:intro-affine-determinant} is nonzero
\citep{Tao2016}.  We record the corresponding log-weighted corollary
and a finite-complement reassembly.  We also prove that the absolute
value of the \(r=s=1\) fiber has positive squarefree density and mass
of order \(J\log^2J\).  Hence a coefficientwise absolute-value estimate
cannot close that component.

The organization is as follows.  \Cref{sec:interface} fixes the exact
primitive packet.  \Cref{sec:dynamic-mean} extends the Ramanujan-mean
bound and removes the scalar ultra-long tail.
\Cref{sec:dynamic-conductor,sec:exponent-strip} prove the dynamic-cutoff
closure and its exponent ledgers.  \Cref{sec:complement-fibers} derives
the complementary affine forms.  \Cref{sec:logarithmic-bridge} gives
the fixed-fiber logarithmic theorem and the absolute-value no-go.
\Cref{sec:barriers-roadmap} states the remaining uniform gate and strict
scope.  The exact certificate and conclusion are in
\cref{sec:certificate-conclusion}.
